{
\begin{table}[!ht]
\fontsize{11pt}{11pt}\selectfont
\setstretch{1}

\begin{tabular}{llll}
\hline

\multicolumn{4}{l}{\textbf{VAASAN YLIOPISTO}} \\ \\
\multicolumn{4}{l}{School of Technology and Innovation} \\ \\

Author: &
\multicolumn{3}{l}{Nahid Hasan Sarkar} \\ \\

Thesis title: &
\multicolumn{3}{l}{\textbf{Developing a Forecasting-Based Decision-Support Framework for Third-Party Delivery Sales Optimization}} \\ \\

Degree: &
\multicolumn{3}{l}{Master of Science in Technology} \\ \\

Degree Specialization: &
\multicolumn{3}{l}{Industrial Systems Analytics} \\ \\

Supervisor: &
\multicolumn{3}{l}{Jyri Naarmala} \\ \\

Year of graduation: &
2026 &
Number of pages: &
\pageref{LastPage} \\                                                               \\

\hline
\\
\multicolumn{4}{l}{\textbf{ABSTRACT:}}                                               \\ \\

\multicolumn{4}{p{14.5cm}}{

This study develops a data-driven decision-support framework to optimize third-party delivery platform sales in the Quick Service Restaurant (QSR) industry. With the rapid growth of platforms such as Wolt and Foodora, large volumes of transactional data are generated, offering opportunities for advanced analytics. However, many businesses do not fully utilize this data for forecasting and decision-making.

The research follows a Design Science Research (DSR) approach to design and evaluate a framework that integrates data preprocessing, demand forecasting, and profit optimization. Historical sales data from Wolt and Foodora are first cleaned and integrated into a unified dataset. A Prophet time series model is then applied to forecast platform-level demand, capturing trend and weekly seasonal patterns. The forecasting results are further used as input for an optimization model that generates profit-oriented decision recommendations under practical constraints such as delivery capacity and platform fees.

The proposed framework is evaluated through simulation and back-testing using historical data. Forecasting performance is assessed using standard accuracy metrics including MAE, RMSE, and MAPE. The results show that the proposed approach can improve decision-making by providing more accurate demand predictions and more efficient platform allocation strategies. The integration of predictive and prescriptive analytics demonstrates clear potential for improving profitability and operational efficiency.

Overall, the study contributes to the field of business analytics by combining predictive and prescriptive analytics into a unified decision-support framework for multi-platform food delivery operations.

} \\

\hline

\multicolumn{4}{p{14.5cm}}{
\textbf{Keywords:} Demand forecasting, Prophet model, decision support systems, optimization, food delivery platforms
} 


\end{tabular}
\end{table}
}